\chapter{生产性粉尘与职业性肺部疾病}
\setcounter{chapter}{4}
\chapterintro{
本章为课程核心章节（15学时，占比最大）。

\textbf{【教学大纲要求】}

\imp{掌握：}生产性粉尘与尘肺的定义；生产性粉尘的来源与分类；粉尘的理化特性及其卫生学意义\keyword{*}；粉尘对健康的影响；粉尘危害的控制（防尘八字方针）；矽肺的定义、影响发病的主要因素、病理变化（病理形态的分型\keyword{*}）、发病机制（主要学说）、临床表现与并发症、X线诊断分期、患者处理；硅酸盐尘肺的共同特点；石棉的理化特性、接触机会、影响发病因素、病理改变与发病机制、临床表现与X线胸片特征\keyword{*}、常见并发症；煤工尘肺（病理改变、临床表现与诊断处理、类风湿性尘肺结节）；职业性变态反应性肺泡炎；棉尘病。

\imp{熟悉：}其他矽酸盐尘肺（滑石肺、云母肺、水泥尘肺）；其他粉尘所致尘肺（电焊工尘肺、铸工尘肺等）；有机粉尘的来源、分类及对健康的危害。

（注：\keyword{*} 标示教学难点）
}

% ----- 4.1 概述 -----
\section{生产性粉尘概述}

\begin{defbox}
\textbf{生产性粉尘}：在生产活动中产生的、能够较长时间悬浮于生产环境中的固体微粒。粒径多为0.1--10 $\mu$m。

\textbf{尘肺（Pneumoconiosis）}：在职业活动中长期吸入生产性矿物性粉尘并在肺内潴留而引起的以肺组织\keyword{弥漫性纤维化}为主的全身性疾病。
\end{defbox}

\subsection{生产性粉尘的来源与分类}

\begin{keybox}
\textbf{来源：}
\begin{enumerate}[leftmargin=*]
    \item 固体物质的机械加工或粉碎（采矿、凿岩、破碎、研磨）
    \item 物质加热时产生的蒸气在空气中冷凝或氧化（熔炼、焊接）
    \item 有机物质不完全燃烧（烟尘）
    \item 粉末状物质的混合、过筛、包装、搬运
\end{enumerate}

\textbf{分类（按性质）：}
\begin{enumerate}[leftmargin=*]
    \item \imp{无机粉尘}：
    \begin{itemize}
        \item 矿物性粉尘：矽尘（石英尘）、石棉尘、煤尘、滑石尘等
        \item 金属性粉尘：铅尘、铁尘、锡尘等
        \item 人工无机粉尘：水泥尘、玻璃纤维等
    \end{itemize}
    \item \imp{有机粉尘}：
    \begin{itemize}
        \item 植物性粉尘：棉尘、木尘、谷物粉尘等
        \item 动物性粉尘：皮毛尘、骨粉等
        \item 人工有机粉尘：塑料尘、染料尘等
    \end{itemize}
    \item \imp{混合性粉尘}：上述两类或两类以上粉尘混合存在
\end{enumerate}
\end{keybox}

\subsection{粉尘的理化性质及其卫生学意义*}

\begin{keybox}
\textbf{1. 粉尘的化学组成}：决定其对机体危害性质的最重要因素。含游离SiO$_2$的粉尘致矽肺；石棉尘致石棉肺；含铅粉尘致铅中毒。

\textbf{2. 粉尘浓度}：单位体积空气中的粉尘重量（mg/m$^3$）。浓度越高，危害越大。

\textbf{3. 粉尘分散度*}：粉尘颗粒大小的组成，用粒径分布的百分数表示。
\begin{itemize}[leftmargin=*]
    \item 粒径越小，悬浮时间越长，被吸入机会越多
    \item \keyword{空气动力学直径（AED）}：不考虑尘粒的形状、大小和密度，如果其在空气中的沉降速度与一种密度为1的球形粒子的沉降速度一样时，则这种球形粒子的直径即为该种粉尘粒子的AED。AED($\mu$m) = $d_p\sqrt{Q}$（$d_p$：光镜下投影直径，Q：粉尘比重）
    \item \keyword{可吸入性粉尘（IP）}：AED 10--15$\mu$m，能进入呼吸道
    \item \keyword{呼吸性粉尘（RP）}：AED <5$\mu$m，能到达肺泡，危害最大
    \item AED >15$\mu$m：多被阻留在鼻腔、咽喉，不易吸入
    \item >10$\mu$m：主要沉积在上呼吸道；<5$\mu$m：可达肺泡
    \item 2--5$\mu$m的粉尘致肺纤维化能力最强
\end{itemize}

\textbf{4. 粉尘的溶解度：}有毒粉尘（铅、砷等）溶解度越大，毒性越强；无毒粉尘溶解度越大，越易排出，危害越小。致纤维化粉尘（石英等）溶解度很低但危害持久。

\textbf{5. 粉尘的荷电性：}同种电荷相互排斥，增加悬浮稳定性；荷电粉尘易滞留于呼吸道。

\textbf{6. 粉尘的形状与硬度：}球形沉降快，片状/纤维状沉降慢；硬质粉尘对呼吸道机械损伤较大。

\textbf{7. 粉尘的爆炸性：}某些粉尘（煤尘、面粉）在空气中达到一定浓度时遇火可引起爆炸。

\textbf{8. 粉尘的吸湿性：}吸湿性强的粉尘易在呼吸道沉积。
\end{keybox}

\subsection{人体对生产性粉尘的防御和清除}

\begin{keybox}
人体呼吸道具有完善的多层次防御和清除机制，能将绝大部分吸入的粉尘排出体外，但当粉尘浓度过高或长期持续接触时，防御清除机制可能被破坏，导致粉尘在肺内潴留。

\textbf{一、鼻腔和咽喉部}
\begin{itemize}[leftmargin=*]
    \item 鼻毛和鼻腔结构：阻挡>10$\mu$m的粉尘颗粒
    \item 鼻甲弯曲结构：产生湍流，增加粉尘粒子撞击并沉积于黏膜表面的机会
    \item 喷嚏反射：将粉尘颗粒排出
\end{itemize}

\textbf{二、气管和支气管——黏液-纤毛清除系统（最重要清除机制）}
\begin{itemize}[leftmargin=*]
    \item 纤毛上皮细胞：以每分钟约1200次的频率向上摆动，推动黏液（含粉尘）以每分钟1--3cm的速度向咽喉部移动
    \item 黏液层：杯状细胞和黏液腺分泌的黏液捕获粉尘颗粒（主要捕获2--10$\mu$m的粉尘）
    \item 咳嗽反射：纤毛推动的含粉尘黏液到达咽喉部后被咳出或咽下
\end{itemize}

\textbf{三、肺泡——肺泡巨噬细胞清除}
\begin{itemize}[leftmargin=*]
    \item 到达肺泡（<5$\mu$m，即呼吸性粉尘）的粉尘主要通过\keyword{肺泡巨噬细胞}吞噬清除
    \item 巨噬细胞吞噬粉尘后：
    \begin{itemize}
        \item 大部分携带粉尘的巨噬细胞沿肺泡表面向终末细支气管移动→纤毛-黏液系统排出
        \item 部分进入肺间质淋巴系统→淋巴结
        \item 部分在肺泡内死亡，释放出吞噬的粉尘颗粒→被其他巨噬细胞再次吞噬
    \end{itemize}
    \item \keyword{巨噬细胞吞噬SiO$_2$后发生损伤和死亡}是矽肺发病的起始环节
\end{itemize}

\textbf{四、淋巴系统清除}
\begin{itemize}[leftmargin=*]
    \item 进入肺间质的粉尘颗粒沿淋巴引流进入肺门淋巴结和纵隔淋巴结
    \item 粉尘在淋巴结中蓄积可引起淋巴结病变（淋巴结肿大、纤维化）
\end{itemize}

\textbf{五、影响清除效率的因素：}
\begin{itemize}[leftmargin=*]
    \item 粉尘浓度和粒径（浓度越高、粒径越小，越易穿透防御系统）
    \item 粉尘的细胞毒性（SiO$_2$对巨噬细胞有毒害作用→抑制清除功能）
    \item 吸烟（抑制纤毛运动和巨噬细胞功能，显著降低清除效率）
    \item 个体因素（年龄、呼吸系统疾病等）
\end{itemize}

\textbf{六、清除机制的局限性：}
\begin{itemize}[leftmargin=*]
    \item 长期高浓度粉尘暴露→清除系统超负荷→粉尘在肺内潴留
    \item 致纤维化粉尘（SiO$_2$、石棉）可损伤清除系统本身（巨噬细胞死亡、纤毛受损）
    \item 潴留粉尘在肺内逐渐形成纤维化病灶
\end{itemize}
\end{keybox}

\subsection{粉尘对健康的影响}

\begin{enumerate}[leftmargin=*]
    \item \imp{致纤维化作用}：引起尘肺病（矽肺、石棉肺、煤工尘肺等）——最严重
    \item \imp{局部刺激作用}：呼吸道炎症（鼻炎、咽炎、支气管炎）
    \item \imp{中毒作用}：吸入铅、锰、砷等有毒粉尘可致全身性中毒
    \item \imp{致癌作用}：石棉致肺癌和间皮瘤（IARC 1类）；放射性粉尘致肺癌
    \item \imp{致敏作用}：有机粉尘可引起变态反应性肺泡炎（如农民肺、棉尘症）
    \item \imp{感染作用}：带有病原菌的粉尘（炭疽杆菌等）
    \item \imp{皮肤损害}：堵塞皮脂腺→毛囊炎、粉刺
    \item \imp{对眼和上呼吸道的损害}：结膜炎、角膜炎
\end{enumerate}

\subsection{我国法定尘肺病（13种）}

\begin{enumerate}[leftmargin=*]
    \item 矽肺、煤工尘肺、石墨尘肺、碳黑尘肺、石棉肺、滑石尘肺、水泥尘肺、云母尘肺、陶工尘肺、铝尘肺、电焊工尘肺、铸工尘肺、其他尘肺（按病因分类的13种）
\end{enumerate}

\subsection{我国尘肺病发病形势}

\begin{itemize}[leftmargin=*]
    \item 截至2023年底，累计报告职业性尘肺病超93.10万例，实际发病数远超报告数，现患约45万例
    \item 矽肺和煤工尘肺占职业病病例报告总数的88.68\%
    \item 特点：尘肺病例占总报告病例的80\%以上；发病工龄缩短；速发型尘肺以煤矿开采和洗选业为主，高发工种为凿岩工、采煤工
\end{itemize}

\subsection{粉尘危害的控制}

\begin{keybox}
\textbf{防尘八字方针：}\keyword{革、水、密、风、护、管、教、查}

\begin{enumerate}[leftmargin=*]
    \item \textbf{革}——技术革新，改革工艺设备和生产方法
    \item \textbf{水}——湿式作业（湿式凿岩、水磨）
    \item \textbf{密}——密闭尘源（密闭输送、密闭粉碎）
    \item \textbf{风}——通风除尘（局部排风、全面通风）
    \item \textbf{护}——个人防护（防尘口罩、防尘服）
    \item \textbf{管}——维护管理（设备维护、制度管理）
    \item \textbf{教}——宣传教育（防尘知识教育）
    \item \textbf{查}——定期检查（环境监测、健康检查）
\end{enumerate}
\end{keybox}

% ----- 4.2 矽肺 -----
\section{矽肺（Silicosis）}

\begin{defbox}
\textbf{矽肺（Silicosis）}：在生产过程中长期吸入含有游离二氧化硅（SiO$_2$）粉尘而引起的以\keyword{肺组织弥漫性纤维化}为主的全身性疾病。是尘肺中\keyword{危害最严重、进展最快、最常见}的一种。
\end{defbox}

\subsection{影响矽肺发病的主要因素}

\begin{enumerate}[leftmargin=*]
    \item 粉尘中\keyword{游离SiO$_2$含量}：含量越高，危害越大，发病越快
    \item 粉尘浓度：浓度越高，发病越早
    \item 粉尘分散度：粒径越小，危害越大
    \item 接触工龄：一般10--15年以上，少数1--2年即可发病（\keyword{快型矽肺}）
    \item 个体因素：年龄、健康状况、遗传因素
    \item 混合性粉尘的联合作用
\end{enumerate}

\subsection{矽肺的病理改变*}

\begin{keybox}
\textbf{基本病理改变：}
\begin{enumerate}[leftmargin=*]
    \item \imp{矽结节}——矽肺的特征性病理改变
    \begin{itemize}
        \item 细胞性结节：巨噬细胞聚集，吞噬SiO$_2$颗粒
        \item 纤维性结节：成纤维细胞增生、胶原纤维形成
        \item 玻璃样结节：胶原纤维透明变性（典型矽结节）
    \end{itemize}
    \item \imp{弥漫性间质纤维化}
    \item \imp{团块状纤维化}——矽结节融合成大块纤维化病变
\end{enumerate}

\textbf{病理形态分裂：}
\begin{itemize}[leftmargin=*]
    \item \imp{结节型}：以矽结节形成为主，多见于接触含矽量高的粉尘
    \item \imp{弥漫纤维化型}：以弥漫性间质纤维化为主，多见于接触含矽量较低的粉尘
    \item \imp{矽性蛋白沉积型}：肺泡腔内含大量蛋白性液体
\end{itemize}
\end{keybox}

\subsection{矽肺的发病机制——主要学说*}

\begin{enumerate}[leftmargin=*]
    \item \imp{机械刺激学说}：早期理论，SiO$_2$颗粒机械刺激引起纤维化
    \item \imp{化学中毒学说（溶解学说）}：SiO$_2$在体液中缓慢溶解形成矽酸（H$_2$SiO$_3$），对细胞产生毒性作用
    \item \imp{免疫学说}：SiO$_2$使巨噬细胞释放抗原物质，引发自身免疫反应
    \item \imp{细胞毒学说}：SiO$_2$损伤肺泡巨噬细胞→释放致纤维化因子→成纤维细胞增生→胶原纤维形成（目前公认的主要机制）：
    \begin{itemize}
        \item SiO$_2$+巨噬细胞→损伤、死亡→释放IL-1、TNF-$\alpha$、FN等→成纤维细胞活化→胶原纤维增生
    \end{itemize}
\end{enumerate}

\subsection{矽肺的临床表现与X线诊断}

\begin{keybox}
\textbf{临床表现：}
\begin{itemize}[leftmargin=*]
    \item 早期可无症状或症状轻微
    \item 主要症状：胸闷、胸痛、气短、咳嗽、咳痰
    \item 体征：早期常无阳性体征；晚期可出现肺气肿、肺心病体征
    \item 并发症：肺结核（\keyword{矽肺结核}，最常见最严重并发症）、慢性阻塞性肺疾病（COPD）、肺部感染、自发性气胸、肺心病、呼吸衰竭
\end{itemize}

\textbf{X线胸片诊断（尘肺诊断的主要依据）*：}
\begin{itemize}[leftmargin=*]
    \item 小阴影（<10mm）：圆形小阴影（p/q/r型）、不规则小阴影（s/t/u型）
    \item 大阴影（$\geq$10mm）：矽结节融合成团块
    \item 胸膜改变
\end{itemize}

\textbf{X线诊断分期（依据GBZ70-2015）：}
\begin{enumerate}[leftmargin=*]
    \item \imp{壹期尘肺（I）}：①总体密集度1级的小阴影，分布范围$\geq$2个肺区；②石棉粉尘接触者，1级小阴影仅1个肺区，伴胸膜斑；③石棉粉尘接触者，小阴影总体密集度为0，$\geq$2个肺区0/1级密集度，伴胸膜斑
    \item \imp{贰期尘肺（II）}：①总体密集度2级的小阴影，分布范围>4个肺区；②总体密集度3级的小阴影，分布范围$\geq$4个肺区；③石棉粉尘接触者，1级小阴影>4个肺区，伴胸膜斑累及心缘/膈面；④石棉粉尘接触者，2级小阴影$\geq$4个肺区，伴胸膜斑累及心缘/膈面
    \item \imp{叁期尘肺（III）}：①大阴影长径$\geq$20mm、短径>10mm；②总体密集度3级的小阴影，分布范围>4个肺区，伴小阴影聚集；③总体密集度3级的小阴影，分布范围>4个肺区，伴大阴影；④石棉粉尘接触者，3级小阴影>4个肺区，伴胸膜斑长度之和>单侧胸壁1/2或累及心缘
\end{enumerate}
\end{keybox}

\subsection{矽肺的处理}

\begin{enumerate}[leftmargin=*]
    \item \imp{调离粉尘作业}：一旦确诊，立即调离
    \item \imp{综合治疗}：药物治疗（克矽平等）、对症治疗、康复治疗
    \item \imp{职业伤残程度鉴定}：按GB/T 16180进行劳动能力鉴定
\end{enumerate}

% ----- 4.3 硅酸盐尘肺 -----
\section{硅酸盐尘肺}

\subsection{概述}

\begin{defbox}
\textbf{硅酸盐（Silicate）}：由SiO$_2$、金属氧化物和结晶水组成的矿物。石棉、滑石、云母、水泥等均为硅酸盐。

\textbf{硅酸盐尘肺的共同特点*：}
\begin{enumerate}[leftmargin=*]
    \item 病理改变主要表现为\keyword{弥漫性间质纤维化}（与矽肺的结节型不同）
    \item 组织切片中可见\keyword{石棉小体}（石棉肺）、滑石小体等
    \item X线胸片以不规则小阴影为主
    \item 致纤维化能力一般弱于石英
\end{enumerate}
\end{defbox}

\subsection{石棉肺（Asbestosis）}

\begin{keybox}
\textbf{石棉的理化性质：}具有纤维结构的硅酸盐矿物。分为蛇纹石类（温石棉，占90\%以上）和角闪石类（青石棉、铁石棉等）。耐热、耐酸碱、绝缘、抗拉强度大。

\textbf{接触机会：}石棉矿开采选矿、石棉制品制造（石棉瓦、石棉板、刹车片）、建筑拆除和维修。

\textbf{影响石棉肺发病的因素：}
\begin{itemize}[leftmargin=*]
    \item 石棉种类：角闪石类（青石棉）>蛇纹石类（温石棉）
    \item 纤维长度和直径：长纤维（>5$\mu$m）、细纤维（<0.25$\mu$m）致病力强
    \item 粉尘浓度和接触时间
    \item 个体因素
\end{itemize}

\textbf{病理改变与发病机制*：}
\begin{itemize}[leftmargin=*]
    \item 弥漫性间质纤维化（与矽肺的结节型不同）
    \item \keyword{石棉小体}：石棉纤维外包绕铁蛋白形成的哑铃状或串珠状结构，是石棉肺的\keyword{病理学标志}
    \item 胸膜斑：壁层胸膜的局限性增厚
\end{itemize}

\textbf{临床表现：}
\begin{itemize}[leftmargin=*]
    \item 症状：呼吸困难（最早、最常见）、咳嗽、胸痛
    \item 体征：双肺底吸气末Velcro啰音（爆裂音），杵状指
    \item X线特征*：不规则小阴影为主（s/t/u型），胸膜斑，肺门结构紊乱
\end{itemize}

\textbf{并发症：}
\begin{itemize}[leftmargin=*]
    \item \imp{肺癌}：石棉致肺癌（IARC 1类致癌物），吸烟与石棉暴露有协同致癌作用
    \item \imp{弥漫性恶性间皮瘤}：胸膜或腹膜间皮瘤，与石棉暴露密切相关（青石棉>温石棉）
\end{itemize}

\textbf{其他石棉相关疾病：}石棉所致胸膜斑、石棉所致胸腔积液、石棉所致支气管肺癌
\end{keybox}

\subsection{其他硅酸盐尘肺}

\begin{table}[H]
\centering
\caption{其他硅酸盐尘肺}
\begin{tabularx}{\textwidth}{lX}
\hline
\textbf{类型} & \textbf{特点} \\
\hline
\imp{滑石肺} & 滑石矿开采和加工；病理以间质纤维化为主，可见滑石小体；可伴胸膜改变 \\
\imp{云母肺} & 云母矿开采加工；纤维化程度较矽肺轻 \\
\imp{水泥尘肺} & 水泥制造；混合性粉尘；病程进展缓慢 \\
\hline
\end{tabularx}
\end{table}

% ----- 4.4 煤工尘肺 -----
\section{煤工尘肺（CWP）}

\begin{defbox}
\textbf{煤工尘肺（Coal Worker's Pneumoconiosis, CWP）}：煤矿工人长期吸入煤尘和含矽混合性粉尘所致的尘肺。包括：
\begin{enumerate}[leftmargin=*]
    \item \imp{煤肺}：吸入纯煤尘（采煤工），病理以煤斑和灶周肺气肿为特征
    \item \imp{煤矽肺}：吸入煤尘和矽尘的混合粉尘（掘进工），兼有煤肺和矽肺的病理特征
\end{enumerate}
\end{defbox}

\begin{keybox}
\textbf{病理改变：}
\begin{itemize}[leftmargin=*]
    \item \imp{煤斑（Coal Macule）}：由含煤尘的巨噬细胞聚集而成，周围伴有\keyword{灶周肺气肿}
    \item 煤结节：煤尘+少量胶原纤维
    \item 进行性大块纤维化（PMF）：晚期煤结节融合成大块
\end{itemize}

\textbf{类风湿性尘肺结节（Caplan综合征）：}煤工尘肺合并类风湿关节炎，X线胸片显示圆形结节影，短期内出现并增大。

\textbf{临床表现与诊断：}胸闷、气短、咳嗽、咳黑痰。X线以小阴影为主。
\end{keybox}

% ----- 三种主要尘肺的对比 -----
\subsection{矽肺、石棉肺与煤工尘肺的比较*}

\begin{keybox}
以下三种尘肺的异同点是考试核心内容，需重点掌握其病理特征和X线表现的区别。

\begin{table}[H]
\centering
\caption{三种主要尘肺的全面对比}
\begin{tabularx}{\textwidth}{lXXX}
\hline
\textbf{对比项目} & \textbf{矽肺（Silicosis）} & \textbf{石棉肺（Asbestosis）} & \textbf{煤工尘肺（CWP）} \\
\hline
\imp{致病粉尘} & 游离SiO$_2$粉尘 & 石棉纤维 & 煤尘（煤肺）或煤尘+矽尘（煤矽肺）\\
\hline
\imp{主要行业} & 采矿、隧道、陶瓷、玻璃 & 石棉矿、石棉制品、建筑拆除 & 煤矿开采（采煤工/掘进工）\\
\hline
\imp{病理类型} & \textbf{结节型}纤维化为主 & \textbf{弥漫性间质}纤维化为主 & 煤斑伴\textbf{灶周肺气肿}为主 \\
\hline
\imp{特征性结构} & \textbf{矽结节}（细胞性→纤维性→玻璃样） & \textbf{石棉小体}（哑铃状/串珠状）、\textbf{胸膜斑} & \textbf{煤斑}（煤尘+巨噬细胞聚集）\\
\hline
\imp{X线表现} & 圆形/类圆形小阴影（p/q/r型）为主；三期可见大阴影 & \textbf{不规则小阴影}（s/t/u型）为主；胸膜斑；肺门结构紊乱 & 小阴影（圆形或不规则）；晚期PMF形成大阴影 \\
\hline
\imp{肺功能损害} & 混合性通气功能障碍 & \textbf{限制性通气功能障碍}为主（早期即可出现），肺顺应性↓ & 阻塞性通气功能障碍为主（单纯煤工尘肺）\\
\hline
\imp{临床症状} & 胸闷、胸痛、气短、咳嗽 & \textbf{呼吸困难}（最早最常见）、Velcro啰音、杵状指 & 胸闷、气短、咳嗽、\textbf{咳黑痰} \\
\hline
\imp{主要并发症} & \textbf{肺结核（矽肺结核）}——最常见最严重 & \textbf{肺癌、弥漫性恶性间皮瘤} & 肺结核、慢性支气管炎、肺气肿 \\
\hline
\imp{致癌性} & IARC 1类（1997） & IARC 1类——肺癌和间皮瘤 & 未明确分类 \\
\hline
\imp{病程进展} & 脱离粉尘后病情仍可继续进展 & 脱离粉尘后可缓慢进展 & 脱离粉尘后进展极缓慢或停止 \\
\hline
\imp{特殊类型} & 快型矽肺（1--2年发病） & 胸膜间皮瘤（青石棉风险最高）& Caplan综合征（合并类风湿）\\
\hline
\hline
\end{tabularx}
\end{table}

\textbf{矽结节与煤尘纤维灶（煤斑）的形态学对比：}
\begin{table}[H]
\centering
\caption{矽结节与煤尘纤维灶的形态学对比}
\begin{tabularx}{\textwidth}{lXX}
\hline
\textbf{对比项目} & \textbf{矽结节} & \textbf{煤尘纤维灶（煤斑）} \\
\hline
成分 & 大量胶原纤维+少量粉尘 & 大量煤尘+少量胶原纤维 \\
质地 & 硬韧、玻璃样变 & 柔软 \\
大小 & 较大（可达数mm至数cm） & 较小（1--3mm） \\
网状纤维 & 早期有，晚期减少 & 丰富 \\
胶原纤维 & 大量，致密透明变性 & 少量，疏松 \\
周围肺气肿 & 结节周围可见 & 明显灶周肺气肿（特征性）\\
可逆性 & 不可逆 & 部分可逆（脱离粉尘后）\\
\hline
\end{tabularx}
\end{table}
\end{keybox}

% ----- 4.5 其他粉尘所致尘肺 -----
\section{其他粉尘所致尘肺}

\begin{table}[H]
\centering
\caption{其他类型尘肺}
\begin{tabularx}{\textwidth}{lX}
\hline
\textbf{类型} & \textbf{致病粉尘及特点} \\
\hline
\imp{电焊工尘肺} & 电焊烟尘（氧化铁、MnO等混合粉尘）；X线以不规则小阴影为主 \\
\imp{铸工尘肺} & 铸造用型砂（含SiO$_2$）；多为轻型尘肺，进展缓慢 \\
\imp{碳黑尘肺} & 碳黑生产和使用 \\
\imp{石墨尘肺} & 天然石墨矿开采；含矽石墨致纤维化较重 \\
\imp{铝尘肺（铝土肺）} & 金属铝粉吸入 \\
\hline
\end{tabularx}
\end{table}

% ----- 4.6 有机粉尘 -----
\section{有机粉尘所致的肺部疾病}

\begin{keybox}
\subsection{职业性变态反应性肺泡炎（EAA）}

\begin{defbox}
\textbf{职业性变态反应性肺泡炎（Extrinsic Allergic Alveolitis）}：反复吸入某些有机粉尘中具有抗原性的物质所引起的以肺泡壁和间质过敏性炎症为主的免疫性肺部疾病。
\end{defbox}

\textbf{常见类型：}
\begin{itemize}[leftmargin=*]
    \item \imp{农民肺}：吸入发霉干草中的嗜热放线菌孢子
    \item \imp{蔗渣肺}：吸入发霉甘蔗渣中的嗜热放线菌
    \item \imp{蘑菇工肺}：吸入蘑菇堆肥中微生物孢子
    \item \imp{饲鸟者肺}：吸入鸟类羽毛和排泄物中的蛋白抗原
\end{itemize}

\textbf{发病机制：}III型（免疫复合物型）和IV型（细胞免疫型）变态反应。

\textbf{临床分期：}急性期（接触后4--8小时发热、咳嗽、呼吸困难）；亚急性期；慢性期（不可逆性肺纤维化）。

\subsection{棉尘症（Byssinosis）}

又称"星期一热"，棉纺工人吸入棉尘后出现的胸部紧束感、呼吸困难，休息日后第一天工作症状最重。机制与棉尘中内毒素有关。
\end{keybox}

% ----- 4.7 复习题 -----
\section{复习题}

\begin{enumerate}[leftmargin=*, itemsep=8pt]
    \item \textbf{【名词解释】}生产性粉尘；尘肺；矽肺；呼吸性粉尘
    \item \textbf{【名词解释】}矽结节；石棉小体；胸膜斑；Caplan综合征
    \item \textbf{【名词解释】}农民肺；棉尘症（星期一热）
    \item \textbf{【简答题】}简述粉尘的理化性质及其卫生学意义。
    \item \textbf{【简答题】}简述人体对生产性粉尘的防御和清除机制。
    \item \textbf{【简答题】}简述防尘八字方针及其内涵。
    \item \textbf{【简答题】}简述矽肺的基本病理改变和发病机制（主要学说）。
    \item \textbf{【简答题】}简述矽肺的X线诊断分期标准。
    \item \textbf{【简答题】}简述硅酸盐尘肺的共同特点及石棉肺的特征性改变。
    \item \textbf{【简答题】}比较矽肺与石棉肺的病理和临床特征。
    \item \textbf{【简答题】}简述煤工尘肺的病理类型和特征。
    \item \textbf{【简答题】}简述有机粉尘所致肺部疾病的类型及特点。
\end{enumerate}

\subsection*{参考答案要点}

\begin{enumerate}[leftmargin=*, itemsep=5pt]
    \item \textbf{生产性粉尘}：生产活动中产生的能较长时间悬浮的固体微粒。\textbf{尘肺}：长期吸入矿物性粉尘所致以肺弥漫性纤维化为主的疾病。\textbf{矽肺}：长期吸入游离SiO$_2$粉尘所致肺弥漫性纤维化，危害最严重。\textbf{呼吸性粉尘}：粒径$\leq$5$\mu$m，能到达肺泡的粉尘。
    \item \textbf{矽结节}：矽肺特征性病理改变，细胞性→纤维性→玻璃样结节。\textbf{石棉小体}：石棉纤维外包铁蛋白形成的哑铃状结构，石棉肺病理学标志。\textbf{胸膜斑}：壁层胸膜局限性增厚。\textbf{Caplan综合征}：煤工尘肺合并类风湿关节炎的圆形结节影。
    \item \textbf{农民肺}：吸入发霉干草中嗜热放线菌孢子所致变态反应性肺泡炎。\textbf{棉尘症}：棉纺工人吸入棉尘出现的胸部紧束感，休息日后第一天最重（星期一热）。
    \item 八项：化学组成（决定危害性质）、浓度、分散度（粒径越小危害越大，呼吸性粉尘$\leq$5$\mu$m）、溶解度、荷电性、形状硬度、爆炸性、吸湿性。
    \item 六层防御清除机制：鼻腔（阻挡>10$\mu$m，鼻甲产生湍流）→气管支气管（黏液-纤毛清除系统，最重要，捕获2--10$\mu$m粉尘）→肺泡（巨噬细胞吞噬<5$\mu$m呼吸性粉尘，SiO$_2$损伤巨噬细胞是矽肺起始环节）→淋巴系统（进入肺门和纵隔淋巴结）→清除受限因素（浓度过高、细胞毒性粉尘、吸烟抑制清除）。最终潴留粉尘在肺内引起纤维化。
    \item 革（技术革新）、水（湿式作业）、密（密闭尘源）、风（通风除尘）、护（个人防护）、管（维护管理）、教（宣传教育）、查（定期检查）。
    \item 基本病理：矽结节（细胞性→纤维性→玻璃样）、弥漫性间质纤维化、团块状纤维化。主要学说：细胞毒学说（公认）——SiO$_2$损伤巨噬细胞→释放致纤维化因子→成纤维细胞增生→胶原纤维形成。
    \item 壹期：总体密集度1级小阴影，分布≥2个肺区；或石棉接触者1级1个肺区伴胸膜斑。贰期：密集度2级小阴影分布>4个肺区；或3级分布≥4个肺区；或石棉接触者相应条件伴胸膜斑。叁期：大阴影长径≥20mm、短径>10mm；或3级小阴影分布>4个肺区伴小阴影聚集/大阴影。
    \item 共同特点：弥漫性间质纤维化为主；可见特征性小体；不规则小阴影为主。石棉肺特征：石棉小体（病理标志）、胸膜斑、Velcro啰音、肺癌和间皮瘤。
    \item 矽肺：结节型纤维化、矽结节、圆形小阴影、并发症以肺结核为最常见。石棉肺：弥漫性纤维化、石棉小体、不规则小阴影、胸膜斑、并发症以肺癌和间皮瘤为主。
    \item 煤肺（纯煤尘，煤斑和灶周肺气肿）和煤矽肺（混合粉尘，兼有煤肺和矽肺特征）。晚期可出现进行性大块纤维化（PMF）。
    \item 职业性变态反应性肺泡炎（农民肺、蔗渣肺、蘑菇工肺——III/IV型变态反应）；棉尘症（星期一热——内毒素机制）。
\end{enumerate}

\newpage

% =====================================================================
%  CHAPTER 5: 物理因素及其对健康的影响
% =====================================================================
\newpage